%%
%% fall-lecture12.tex
%% 
%% Made by Alex Nelson <pqnelson@gmail.com>
%% Login   <alex@lisp>
%% 
%% Started on  2026-10-24T10:03:59-0700
%% Last update 2026-10-24T10:03:59-0700
%% 

\lecture{}

\subsection{Vorticity}

\begin{definition}
The \define{Vorticity} for a fluid flow is the [pseudo]vector field $\vorticity:=\curl{\vec{u}}$.
\end{definition}


\begin{definition}
A flow is called \define{Irrotational} if its vorticity vanishes $\vorticity=\vec{0}$.
\end{definition}

\begin{definition}
A vector field $\vec{u}$ is called \define{Solenoidal} if its
divergence vanishes $\div{\vec{u}}=0$ everywhere.
\end{definition}

\begin{recall}
For a solenoidal vector field $\vec{u}$, there exists a unique vector
field $\vec{A}$ such that $\vec{u}=\curl{\vec{A}}$ and
$\div{\vec{A}}=0$ everywhere. (The $\div{\vec{A}}=0$ condition is
needed for uniqueness of $\vec{A}$, otherwise
$\vec{A}+\gradient{\varphi}$ also works for any arbitrary scalar field $\varphi$.)
\end{recall}

\begin{node}
For any finite continuous vector field $\vec{u}$ such that $\vec{u}$
vanishes as $\|\vec{x}\|\to\infty$, we can write
\begin{equation}
\vec{u}=\gradient{\varphi}+\curl{\vec{a}}
\end{equation}
as the sum of the curl-free portion and the divergence-free
portion. Here $\vec{u}$ being ``continuous'' rules out shock waves.
\end{node}

\subsection{2-Dimensional Streamfunction}

\begin{definition}
We say a scalar function $\psi$ is a \define{Streamfunction} (or
``stream function'') if we can write the velocity vector field as
$\vec{u}=\gradient{\psi}$. In this case, the vector field
$\vec{a}=(0,0,\psi)$ satisfies $\curl{\vec{a}}=\vec{u}$.

In particular, this means $\vec{u}=(u,v,0)$ has $u=\partial_{y}\psi$
and $v=-\partial_{x}\psi$.
\end{definition}

\begin{node}
The incompressible 2-dimensional flow has the continuity equation
\begin{equation}
\partial_{x}u+\partial_{y}v=0.
\end{equation}
This means
\begin{equation}
\partial_{x}(\partial_{y}\psi)+\partial_{y}(-\partial_{x}\psi)=0
\end{equation}
identically.
\end{node}

\begin{node}
If we take the contours where $\psi$ equals a constant, then these
describe stream lines. Really?
Well, recall~(\S\ref{defn:stream-line}), stream lines satisfy
\begin{equation}
\frac{\D x}{u}=\frac{\D y}{v},
\end{equation}
which means
\begin{equation}
u\,\D y-v\,\D x = 0.
\end{equation}
We can plug in the expressions $u=\partial_{y}\psi$ and $v=-\partial_{x}\psi$
to obtain
\begin{equation}
\partial_{y}\psi\,\D y+\partial_{x}\psi\,\D x=\D\psi=0.
\end{equation}
This happens when $\psi$ is constant.
\end{node}

\begin{node}
In 2-dimensions, the vorticity can be written as (overloading the
``$\omega$'' symbol):
\begin{equation}
\vorticity=(0,0,\omega).
\end{equation}
Then we can express the only nonzero component of vorticity using the
streamfunction as:
\begin{subequations}
  \begin{align}
\omega &= \partial_{x}v-\partial_{y}u\\
&=\partial_{x}(-\partial_{x}\psi)-\partial_{y}(\partial_{y}\psi)\\
&=-\nabla^{2}\psi.
  \end{align}
\end{subequations}
\end{node}

\begin{node}
Now, let us examine the Navier--Stokes equations neglecting external
force:
\begin{equation}
\partial_{t}u_{i}+u_{j}\partial_{j}u_{i}=\frac{-1}{\rho}\partial_{i}p+\nu\partial_{j}^{2}u_{i}.
\end{equation}
Taking the curl of both sides gives us
\begin{equation}
\partial_{t}\omega_{i}+u_{j}\partial_{j}\omega_{i}-\omega_{j}\partial_{j}u_{i}=0+\nu\partial_{j}^{2}\omega_{i}.
\end{equation}
Usually this is rewritten as
\begin{equation}
\partial_{t}\omega_{i}+u_{j}\partial_{j}\omega_{i}=\omega_{j}\partial_{j}u_{i}+\nu\partial_{j}^{2}\omega_{i},
\end{equation}
and we interpret the first term on the right-hand side as describing
``vortex stretching'', and the second term as the ``viscuous diffusion
of vorticity''.
\end{node}

\begin{node}
Specifically, in 2-dimensions, $\vorticity$ has 1 nonzero component
and $\vec{u}$ has two nonzero components, so
\begin{equation}
\omega_{j}\partial_{j}u_{i}=0
\end{equation}
identically. So the vortex stretching term vanishes in 2-dimensional fluids.
\end{node}

\begin{node}
Then the Navier--Stokes equations in 2-dimensions can be written using
the streamfunction $\psi$ and nonzero vorticity component $\omega$ as
the two equations:
\begin{subequations}
  \begin{align}
\partial_{t}\omega+\partial_{y}\psi\partial_{x}\omega-\partial_{x}\psi\partial_{y}\omega=\nu\partial_{j}^{2}\omega\\
\intertext{and}
-\nabla^{2}\psi=\omega.
  \end{align}
\end{subequations}
These equations are called the \define{Vorticity--Streamfunction Formulation}
of the Navier--Stokes equations in 2-dimensions for incompressible fluids.
\end{node}

\begin{remark}
In three-dimensions, people do not tend to use the
vorticity--streamfunction formulation, it becomes as complicated as
the ``vanilla'' Navier--Stokes equations.
\end{remark}

\begin{remark}[History]
Lagrange's \textit{M\'{e}moire sur la th\'{e}orie du mouvement des fluides}
(1781) introduced Streamfunctions for 2-dimensional incompressible
fluids. Stokes later~\cite{stokes1842steady} extended streamfunctions
to three-dimensional incompressible fluids.
\end{remark}

\subsection{Dynamic Similarity}

\begin{node}
There are two ways to think of nondimensionalization. So far, all
equations use dimensionful quantities like density $\rho$ (dimensions
$\mathsf{M}\cdot\mathsf{L}^{-3}$), and so on. There are two ways to
get nondimensional quantities.
\end{node}

\begin{node}[Formal scaling of equations of motion]
We have variables with length scale, velocity scale,and pressure. We
use stars to indicate dimensionful quantities like $x^{*}$ (dim:
$\mathsf{L}$), $u^{*}$ (dim: $\mathsf{L}/\mathsf{T}$), $t^{*}$.
We can construct nondimensional quantities if we have a length scale
$L$ and a time scale $T$ (or velocity scale $U$) by
\begin{subequations}
  \begin{align}
x &= \frac{x^{*}}{L}\\
u &= \frac{u^{*}}{U}\\
t &= \frac{t^{*}}{U/L},
  \end{align}
\end{subequations}
and so on, where now $u$, $t$, $x$ are dimensionless quantities.
\end{node}

\begin{node}[Continuity equation]
The continuity equation for incompressible fluids are
\begin{equation}
\frac{\partial u^{*}_{i}}{\partial x^{*}_{i}}=\partial^{*}_{i}u^{*}_{i}=0.
\end{equation}
We plug in the $u^{*}_{i}=u_{i}U$ and $x^{*}_{i}=x_{i}L$ to find
\begin{equation}
\frac{\partial(u_{i}U)}{\partial(x_{i}L)} = \frac{U}{L}\partial_{i}u_{i}=0.
\end{equation}
We can ``pull out'' the constants $U$ and $L$ from the differentiation
(this is a bit handwavy but we can make it rigorous). Since $U$ and
$L$ are nonzero quantities, we can divide both sides by $U/L$ to find
\begin{equation}
\partial_{i}u_{i}=0
\end{equation}
is the continuity equation with dimensionless quantities.
\end{node}

\begin{node}[Momentum equation]
The conservation of linear momentum equations can be written down
using dimensionful quantities as:
\begin{equation}
\partial_{t}^{*}u_{i}^{*}+u_{j}^{*}\partial^{*}_{j}u^{*}_{i}=\frac{-1}{\rho}\partial^{*}_{i}p^{*}+\nu(\partial_{j}^{*})^{2}u^{*}_{i}.
\end{equation}
The density $\rho$ is not really an ``unknown'' in the incompressible
fluid, which is why we omit the ``star'' for it. Now, we can rewrite
this using dimensionless quantities as
\begin{equation}
\frac{\partial(u_{i}U)}{\partial(tL/U)}+(u_{j}U)\frac{\partial(u_{i}U)}{\partial(x_{j}L)}=\frac{-1}{\rho}\frac{\partial(p\rho
U^{2})}{\partial(x_{i}L)} + \nu\frac{\partial^{2}(u_{i}U)}{\partial(x_{j}L)\partial(x_{j}L)}.
\end{equation}
Now we can ``pull out the constants'' $L$, $U$ to obtain:
\begin{equation}
\frac{U^{2}}{L}(\partial_{t}u_{i}+u_{j}\partial_{j}u_{i})=\frac{U^{2}}{L}[-\partial_{i}p+\frac{\nu}{UL}\partial_{j}^{2}u_{i}].
\end{equation}
We can divide both sides by $U^{2}/L$ and we obtain the conservation
of linear momentum equations
\begin{equation}
\partial_{t}u_{i}+u_{j}\partial_{j}u_{i}=-\partial_{i}p+\frac{\nu}{UL}\partial_{j}^{2}u_{i}.
\end{equation}
Observe that the only quantity which distinguishes how one
incompressible fluid flows from another is the nondimensional
quantity $\nu/UL$. This number is so important, it merits its own
name: The \define{Reynolds Number} is defined to be
\begin{equation}
\Rey := \frac{UL}{\nu}.
\end{equation}
Importantly: completely different fluids, at completely different
scales, can have identical flows if they have identical Reynolds numbers.
\end{node}

\begin{remark}
The \emph{process} of nondimensionalization we just worked through is
not a unique process (e.g., we might have picked a characteristic time
scale $T$ to our problem instead of a characteristic fluid velocity
scale $U$). We want to pick the process which produces 
cleaner equations of motions.
\end{remark}

\begin{node}[Second approach: Informal Analysis]
The second approach to nondimensionalization is ``informal analysis''
or ``scaling analysis''. This is a much more physical argument. We
look at the Reynolds number as the ratio
\begin{equation}
\Rey = \frac{\mbox{(inertial forces)}}{\mbox{(viscuous forces)}}.
\end{equation}
The inretial forces per unit volume scales like $\rho U^{2}/L$ because
mass scales like $\rho L^{3}$ and acceleration is $U^{2}/L$. Putting
this together,
\begin{equation}
\begin{pmatrix}\mbox{inertial forces}\\
\mbox{per unit volume}
\end{pmatrix}\sim\frac{(\rho L^{3})(U^{2}/L)}{L^{3}}=\frac{\rho U^{2}}{L}.
\end{equation}
The viscuous forces are typically the shear stress, which scales like
$\mu\partial_{y}u\sim\mu U/L$. Then
\begin{equation}
\begin{pmatrix}\mbox{viscuous forces}\\
\mbox{per unit volume}
\end{pmatrix}\sim\frac{(\mu U/L)L^{2}}{L^{3}}=\frac{\mu U}{L^{2}}.
\end{equation}
Then the Reynolds number looks like
\begin{equation}
\Rey\sim\frac{(\rho U^{2}/L)}{(\mu U/L^{2})}=\frac{\rho UL}{\mu}.
\end{equation}
\end{node}

\begin{example}[Drag vs lift]
We have some object of size $L$ moving relative to some fluid at
velocity $U$:
\begin{center}
  \includegraphics{img/img.18}
\end{center}
Then the \define{Drag Coefficient}
\begin{subequations}
  \begin{align}
C_{D} = \frac{\mbox{drag forces}}{\mbox{inertial forces}}=\frac{\mbox{drag}}{\frac{1}{2}\rho U^{2}L^{2}}\label{eq:lec12:me101a:drag-coef}\\
\intertext{and \define{Lift Coefficient}:}
C_{L} = \frac{\mbox{lift forces}}{\mbox{inertial forces}}=\frac{\mbox{lift}}{\frac{1}{2}\rho U^{2}L^{2}}
  \end{align}
\end{subequations}
describe how drag (resp., lift) compare to inertial forces.
\end{example}

\begin{example}
The \define{Froude Number} is defined as
\begin{equation}
\Fr := \sqrt{\frac{\mbox{inertial forces}}{\mbox{external force}}}.
\end{equation}
The square root is a historical artifact, not essential to the physics.
For example, when gravity is the dominant external force,
\begin{equation}
\Fr = \frac{U}{\sqrt{gH}}
\end{equation}
where $H$ is the characteristic length scale (usually, height).
\end{example}

\begin{example}
The \define{Mach Number} is the ratio of the velocity of the fluid to
the velocity of sound propagating through that fluid
\begin{equation}
\Mach = \frac{\mbox{velocity of fluid}}{\mbox{velocity of sound}}=\sqrt{\frac{\mbox{inertial force}}{\mbox{compression force}}}
\end{equation}
\end{example}

\begin{node}
There are also the important \define{Weber Number}
\begin{equation}
\We = \frac{\mbox{inertial force}}{\mbox{surface tension}}
\end{equation}
and the \define{Euler Number}
\begin{equation}
\Eu = \frac{\mbox{Pressure force}}{\mbox{inertial forces}}.
\end{equation}
\end{node}